% ==========================================
% LaTeX CV Reference Variables File
% Fully self-contained configuration
% ==========================================

% --- PERSONAL INFO & HEADER ---
\newcommand{\CVFirstName}{Abderrahmene}
\newcommand{\CVLastName}{KABAR}
\newcommand{\CVJobTitle}{ Consultant Junior Stratégie IT \&  \\[0.15em]   Pilotage de Grands Programmes}
\newcommand{\CVApplicationRef}{test}
\newcommand{\CVPhone}{+33 6 12 34 56 78}
\newcommand{\CVEmail}{abderrahmene.kabar@ecole.ensicaen.fr}
\newcommand{\CVLocation}{Caen, France}
\newcommand{\CVLinkedInURL}{https://www.linkedin.com/in/abderrahmene-kabar}
\newcommand{\CVLinkedInText}{linkedin.com/in/abderrahmene-kabar}
\newcommand{\CVGithubURL}{https://github.com/Impactabdou}
\newcommand{\CVGithubText}{@Impactabdou}
\newcommand{\CVDrivingLicense}{Permis B}
\newcommand{\CVMobility}{Mobile sur toute la France}

% --- PROFESSIONAL PROFILE / SUMMARY ---
\newcommand{\CVProfileText}{À la recherche d’un \textbf{stage de fin d’études de 6 mois à partir de février 2027} au sein de la practice CTO Advisory. Je souhaite accompagner les DSI dans le pilotage de leurs grands programmes de transformation. Mon double cursus ENSICAEN / EM Normandie me permet d’allier la compréhension des architectures IT complexes à l'excellence opérationnelle, financière et organisationnelle (PMO, conduite du changement).}

% --- EDUCATION ---
\newcommand{\EduOneSchool}{EM Normandie}
\newcommand{\EduOneLocation}{Caen, France}
\newcommand{\EduOneDegree}{Programme Grande École (Management \& Stratégie)}
\newcommand{\EduOneDates}{Sept. 2025 -- Août 2027}
\newcommand{\EduOneDesc}{}
\newcommand{\EduOneModules}{\textbf{Modules} : Agilité digitale, Conduite du changement \& Management des organisations, Contrôle budgétaire \& Gestion financière, Diagnostic stratégique \& Business Models, Négociation.}
\newcommand{\EduTwoSchool}{ENSICAEN}
\newcommand{\EduTwoLocation}{Caen, France}
\newcommand{\EduTwoDegree}{Diplôme d'Ingénieur en Informatique (e-Paiement \& Cybersécurité)}
\newcommand{\EduTwoDates}{Sept. 2024 -- Août 2027}
\newcommand{\EduTwoDesc}{\textbf{Résultat actuel :} Mention Bien.}
\newcommand{\EduTwoModules}{}
\newcommand{\EduThreeSchool}{Université de Caen Normandie}
\newcommand{\EduThreeLocation}{Caen, France}
\newcommand{\EduThreeDegree}{Licence en Informatique (Intelligence Artificielle \& Systèmes d'Information)}
\newcommand{\EduThreeDates}{Sept. 2023 -- Août 2024}
\newcommand{\EduThreeDesc}{\textbf{Résultat :} Top 3\% de la promotion (5/160) avec Mention Bien.}

% --- PROFESSIONAL EXPERIENCE ---
\newcommand{\ExpOneCompany}{Groupe Orange}
\newcommand{\ExpOneLocation}{Caen, France}
\newcommand{\ExpOneJob}{Stage Assistant Ingénieriste Sécurisation DNS}
\newcommand{\ExpOneDates}{Avril 2026 -- Août 2026}
\newcommand{\ExpOneTaskA}{Conception et pilotage en méthodologie \textbf{Agile Scrum d'un Proof of Concept (PoC)} répliquant une infrastructure DNS d'Orange afin d'évaluer une solution innovante de sécurisation DNS.}
\newcommand{\ExpOneTaskB}{Modernisation et automatisation à 100\% du cycle de vie des déploiements via une approche \textbf{Infrastructure as Code (IaC)} en utilisant \textbf{Kubernetes, Ansible, Helm et des pipelines CI/CD}.}
\newcommand{\ExpOneTaskC}{Réalisation d'une \textbf{analyse normative des risques} et présentation argumentée des résultats à l'équipe DNS.}
\newcommand{\ExpOneTaskD}{Validation fonctionnelle du PoC et livraison d'une feuille de route stratégique à l'équipe DNS, sécurisant le futur passage en production.}
\newcommand{\ExpTwoCompany}{Chronodrive}
\newcommand{\ExpTwoLocation}{Caen, France}
\newcommand{\ExpTwoJob}{Préparateur et coordinateur de commandes}
\newcommand{\ExpTwoDates}{Fév. 2024 -- Mai 2026}
\newcommand{\ExpTwoTaskA}{Gestion des flux logistiques, coordination des commandes clients et maintien des indicateurs de performance dans un environnement exigeant et rythmé.}

% --- RELEVANT PROJECTS ---
\newcommand{\ProjOneName}{\textbf{Moteur de simulation} $|$ \emph{GitLab CI/CD, Gradle, Agile Kanban}}
\newcommand{\ProjOneDates}{Sept. 2025 -- Déc. 2025}
\newcommand{\ProjOneType}{projet académique}
\newcommand{\ProjOneTaskA}{Développement d'un simulateur au sein d'une équipe de 8 personnes en méthodologie Agile.}
\newcommand{\ProjOneTaskB}{Application de patrons de conception pour structurer le système, couplée à l'automatisation des tests et des déploiements.}
\newcommand{\ProjTwoName}{\textbf{Benchmark d'IA décisionnelles} $|$ \emph{Analyse de données, Modélisation}}
\newcommand{\ProjTwoDates}{Jan. 2024 -- Avr. 2024}
\newcommand{\ProjTwoType}{projet académique}
\newcommand{\ProjTwoTaskA}{Conception et implémentation d'algorithmes d'aide à la décision au sein d'une équipe de 4 étudiants.}
\newcommand{\ProjTwoTaskB}{Collecte, modélisation et analyse de vastes jeux de données afin d'évaluer les performances d'algorithmes et de structurer un benchmark comparatif.}

% --- TECHNICAL SKILLS & LANGUAGES ---
\newcommand{\SkillsLanguages}{Python, Go, Bash, Java (SE/EE), JS, C, C++, SQL, PSQL.}
\newcommand{\SkillsTools}{AWS, Azure, GitLab (CI/CD), Github Actions, Kubernetes, Docker, Ansible, Helm, Helmfile, Prometheus, Grafana.}
\newcommand{\SoftSkills}{Benchmark de solutions, Méthodologies Agiles (Scrum, Kanban, SAFe), Pack Office.}
\newcommand{\SkillsSpoken}{Français (Langue maternelle), Anglais (Courant - \textbf{TOEIC : 910/990}), Arabe (Langue maternelle).}

% --- INTERESTS & HOBBIES ---
\newcommand{\HobbyOne}{\textbf{Expédition \& Randonnée} : Traversée de l'Atlas tellien (Algérie) pendant 2 semaines.}
\newcommand{\HobbyTwo}{\textbf{Esport \& Simulation de gestion} : Pratique compétitive en équipe (Ligue Classée sur Fortnite).}
\newcommand{\HobbyThree}{\textbf{Développement \& Open Source} : Travail préparatoire en cours pour une contribution sur GitHub (Ansible).}

% --- PROFESSIONAL REFERENCES ---
\newcommand{\RefOneName}{M. Mickael LELIARD}
\newcommand{\RefOneRole}{Responsable adjoint département Services Communs d'Infrastructures chez Orange France.}
\newcommand{\RefOneEmail}{mickael.leliard@orange.com}
\newcommand{\RefOnePhone}{+33 6 45 88 36 15}
\newcommand{\RefTwoName}{M. Anthony BLIN}
\newcommand{\RefTwoRole}{Ingénieriste réseaux et infrastructure chez Orange France.}
\newcommand{\RefTwoEmail}{anthony.blin@orange.com}
\newcommand{\RefTwoPhone}{+33 6 48 19 31 01}

